
Consequently, \(E\) is a central extension of
\(U\simeq\boldsymbol\alpha_2\times\boldsymbol\alpha_2\) by
\(\boldsymbol\mu_2\).  Each factor \(\boldsymbol\alpha_2\) of
\(U\) lifts uniquely in \(E\), but \(U\) itself does not lift in
\(E\), since \([X,Y]\ne0\).

\SGASection{7}{Nilpotent Affine Algebraic Groups}
\SGARefTarget{sga3:xvii:section:7}

\SGASubsection{7.1}{Extensions of Groups of Multiplicative Type}
\SGARefTarget{sga3:xvii:subsection:7:1}

\SGAStatement{Proposition 7.1.1}{Let \(S\) be a prescheme, let
\(H\) and \(K\) be two \(S\)-group preschemes of multiplicative
type and finite type, and let \(E\) be an \(S\)-group prescheme
which is an extension of \(K\) by \(H\) (that is, \(K\) is the
quotient of \(E\) by \(H\) for the fpqc topology).  Then \(E\) is
of multiplicative type in either of the following two cases:
\begin{enumerate}
\item[(a)] \SGARefTarget{sga3:xvii:proposition:7:1:1:item:a}
  \(E\) is commutative.
\item[(b)] \SGARefTarget{sga3:xvii:proposition:7:1:1:item:b}
  The fibres of \(K\) are connected.
\end{enumerate}}
\SGARefTarget{sga3:xvii:proposition:7:1:1}

\SGAProof{Proof}
\begin{enumerate}
\item[\SGAInlineRefTarget{sga3:xvii:proposition:7:1:1:item:i}(i)] Suppose first that \(S\) is the spectrum of a field
  \(k\).  The assertion is local for the fpqc topology, so we may
  suppose that \(k\) is algebraically closed and hence that \(H\)
  and \(K\) are diagonalizable.  Notice that \(K\) acts trivially
  on \(H\) by inner automorphisms: this is clear in case \SGARefLink{sga3:xvii:appendix:A}{(a)}, and
  follows from \SGARefLink{sga3:xvii:lemma:6:2:1}{6.2.1} in case \SGARefLink{sga3:xvii:section:b}{(b)}.  Induction on the length of a
  suitable composition series of \(K\) reduces us to the case
  where \(K\) is of one of the following three types:
  \[
    \text{(a) }K=\mathbb G_m,\qquad
    \text{(b) }K=\boldsymbol\mu_p,\qquad
    \text{(c) }K=\boldsymbol\mu_q,\quad(q,p)=1,
  \]
  and in case \SGARefLink{sga3:xvii:proposition:7:1:1:item:c}{(c)}, \(E\) is commutative.  Now embed \(H\) in
  \(\mathbb G_m^r\).  It follows that \(E\) embeds in an extension
  of \(K\) by \(\mathbb G_m^r\), and we may therefore suppose that
  \(H=\mathbb G_m^r\).

  \begin{enumerate}
  \item[(a)] If \(K=\mathbb G_m\), then \(E\) is a smooth,
    connected, affine algebraic group (Exposé VI\textsubscript{B},
    \SGARefLink{prop:sga3-VIB-9.2}{9.2}) of unipotent rank zero, and is therefore a torus.

  \item[(b)] Let \(K=\boldsymbol\mu_p\).  In this case \(E\) is a
    trivial extension.  Indeed, since \(H\) is smooth, the
    canonical morphism
    \[
      \operatorname{Lie}E\longrightarrow
      \operatorname{Lie}\boldsymbol\mu_p
    \]
    is surjective (Appendix II, \SGARefLink{sga3:xvii:proposition:3:2}{3.2}), and the assertion follows
    from \SGARefLink{sga3:xvii:lemma:5-8-2:item:i}{5.8.2(i)}, in view of Appendix II, \SGARefLink{sga3:ii:subsection:2-2}{2.2}.

  \item[\SGAInlineRefTarget{sga3:xvii:proposition:7:1:1:item:c}(c)] Let \(K=\boldsymbol\mu_q\), where \((q,p)=1\), and
    suppose that \(E\) is commutative.  Again the extension \(E\)
    is trivial.  Indeed, let \(x\in E(\bar k)\) lift a generator
    \(\bar x\) of \(\boldsymbol\mu_q(\bar k)\).  The element
    \(x^q\) belongs to \(H(\bar k)\), and hence has the form
    \(y^q\), with \(y\in H(\bar k)\) (notice that
    \(\mathbb G_m^r(\bar k)\) is \(q\)-divisible).  Since \(E\)
    is commutative, \(y^{-1}x\) is a lift of \(\bar x\) having
    order \(q\).
  \end{enumerate}

\item[(ii)] Consider now the general case.  The groups \(H\) and
  \(K\) are flat, affine, and of finite presentation over \(S\)
  (\SGARefLink{sga3:ix:proposition:2:1}{Exposé IX, 2.1}), and consequently so is \(E\) (Exposé
  VI\textsubscript{B}, \SGARefLink{prop:sga3-VIB-9.2}{9.2}).  Using the general technique of
  Exposé VI\textsubscript{B}, \S10, we reduce to the case where
  \(S\) is noetherian.  To prove that \(E\) is of multiplicative
  type, it then suffices to prove that \(E\) is commutative and
  that \({}_nE\) is finite over \(S\) for every \(n\) (Exposé X,
  4.8(b)).

  \begin{enumerate}
  \item[(a)] \(E\) is commutative.  We must verify that the
    morphism
    \[
      E\times_S E\longrightarrow E,\qquad
      (x,y)\longmapsto[x,y]=xyx^{-1}y^{-1}
    \]
    factors through the unit section of \(E\), and it is enough
    to verify this when \(S\) is the spectrum of an artinian
    local ring.  But then \(E\) is of multiplicative type by \SGARefLink{sga3:xvii:proposition:7:1:1:item:i}{(i)}
    and \SGARefLink{sga3:x:corollary:2:3}{Exposé X, 2.3}.
